\section{Related Work}

\subsection{Architectural Erosion and Conformance Checking}

As software systems evolve, changes to their implementation may
cause deviations from the intended architecture, leading to
architectural erosion. Detecting and controlling such deviations
has therefore been an important topic in software architecture
research. Different approaches have been proposed to identify
architectural inconsistencies and maintain alignment between
architectural designs and their implementations.

De Silva and Balasubramaniam \cite{desilva2012}
surveyed approaches for controlling architectural erosion,
covering architecture conformance checking, design enforcement,
evolution management, and architecture restoration. Their work
provides an overview of the techniques available for identifying
and mitigating architectural inconsistencies.

Passos et al. \cite{passos2010} presented an overview of
static architecture-conformance checking techniques, including
dependency-structure matrices, source-code query languages, and
reflexion models. These techniques offer different mechanisms
for detecting structural deviations between an intended
architecture and an implementation. While these studies establish
the foundations of architectural conformance checking, they
primarily focus on the techniques used to represent and detect
architectural inconsistencies.

\subsection{Executable Architecture Constraints and CI Integration}

Architectural constraints are useful only when they can be
checked against software artifacts as the system evolves.
Manually inspecting implementations for architectural violations
can be difficult to sustain, motivating research into executable
constraints and automated conformance checking. Integrating these
checks into development workflows provides another mechanism
for detecting violations during software evolution.

Kallel et al. \cite{kallel2017generating} proposed a process for
generating reusable and executable architecture constraints from
model-based specifications. Their approach transforms design-time
constraints into executable components that can be invoked to
check architectural properties at the implementation level.

From the perspective of evolutionary architecture, Ford et al.
\cite{ford2017} discussed architecture fitness functions
as a means of automatically assessing whether selected
architectural characteristics remain satisfied as a system
evolves. This perspective emphasizes the role of automated
checks in continuously evaluating architectural and other
system properties.

Pinto et al. \cite{pinto2017} investigated the
integration of architectural conformance checking into continuous
integration (CI). They introduced a conformance process in which
architecture checks are triggered during code integration and
violations are addressed before subsequent releases. Their
evaluation in a real-world Java project demonstrated how
executable conformance checks can be incorporated into an
established development workflow. Collectively, these studies
address the generation, specification, and execution of
architectural constraints, as well as their integration into
software development processes.

\subsection{Architectural Conformance in Practice}

Although architectural conformance techniques have been proposed
and evaluated, maintaining consistency between intended and
implemented architectures remains challenging in practice.
The availability of tools does not necessarily ensure their
adoption, nor does detecting an architectural violation guarantee
that it will be addressed. Understanding practitioners' existing
practices and the barriers to adopting conformance techniques is
therefore an important aspect of architecture research.

Ali et al. \cite{ali2018architecture} investigated the state of
practice of architecture consistency through interviews with
19 experienced software engineers. Their study examined the
practices used to maintain architecture consistency, barriers
to adopting explicit approaches, situations in which such
approaches are useful, and practitioners' requirements for
architecture-consistency tools. The findings indicate that
practitioners often rely on informal practices and encounter
barriers to adopting more explicit approaches.

Pinto et al. \cite{pinto2017} also highlighted the
practical challenges of architectural conformance, noting that
existing tools were underused and that detected violations
were rarely corrected. Their proposed CI-based process sought
to address these challenges by making conformance checking
part of the integration workflow. These studies illustrate
that architectural conformance involves not only the technical
capability to detect violations but also the practices through
which checks are incorporated into development and violations
are handled.

Denisov-Blanch et al.~\cite{denisov2026ase} report a related
maintenance pattern for committed AI-configuration artifacts,
such as behavioral rules and tool settings, in 441 corporate
repositories from 27 organizations. At the artifact level,
73.8\% of these configurations are committed once and never
modified. The underlying git histories cover January 2025
through February 2026, about 14 months, and the files are
AI-configuration artifacts rather than executable architecture
rules. Both studies count files that receive no later edit, but
the observation window and the sampled population differ from
the rule-file histories analyzed here. We therefore treat the
73.8\% figure as contextual evidence of set-and-forget
configuration practice, not as a benchmark for our
never-modified share.

% Summary box
\begin{tcolorbox}[
    colback=gray!5,
    colframe=black!50,
    title={Summary},
    fonttitle=\bfseries,
    boxrule=0.5pt,
    arc=1mm,
    left=2mm,
    right=2mm,
    top=1mm,
    bottom=1mm
]
Prior work has examined architectural erosion and conformance
checking techniques, approaches for specifying and executing
architecture constraints, their integration into continuous
integration workflows, and the practical challenges of
maintaining architecture consistency.
A study of committed AI-configuration artifacts finds that
73.8\% are committed once and never modified, within a
14-month corporate sample that we do not use as a benchmark.
However, comparatively little is known about the prevalence
and long-term evolution of executable architecture rules
across large populations of open-source repositories.
\end{tcolorbox}